\documentclass[12pt]{article}
\usepackage[utf8]{inputenc}
\usepackage[spanish]{babel}
\usepackage{amsmath}
\usepackage{amsfonts}
\usepackage{amssymb}
\usepackage{geometry}
\geometry{letterpaper, margin=2.5cm}

\begin{document}

\begin{center}
    \Large \textbf{Evaluación de Examen 1: Unidad I} \\
    \large Física para Ciencias de la Salud (005-1713) \\
    \vspace{0.5cm}
    \textbf{Estudiante:} Rosmeilys Sucre \quad \textbf{C.I.:} 33.386.357
    \vspace{0.3cm} \\
    \large \textbf{Calificación Final: 8/10 puntos}
\end{center}

\vspace{0.5cm}

\section*{Análisis de Resultados}

\subsection*{1. Ángulo plano (Conversión a radianes)}
\textbf{Procedimiento y resultado:} Plantea una regla de tres empleando la equivalencia $90^\circ = \frac{\pi}{2}$ rad. Desarrolla las operaciones adecuadamente obteniendo la fracción $\frac{37,5\pi}{90}$ y concluye con el resultado decimal correcto de $0,41\pi$ Rad. \\
\textbf{Veredicto:} \textbf{Correcto} (2/2 pts).

\subsection*{2. Sistemas de coordenadas (Longitud del hueso)}
\textbf{Procedimiento y resultado:} Calcula correctamente las longitudes de los catetos restando las coordenadas en cada eje ($A=6$ y $B=8$). Al sustituir en la fórmula del Teorema de Pitágoras comete un pequeño error de transcripción escribiendo $7^2$ en lugar de $8^2$. Sin embargo, la operación subsiguiente es correcta ($\sqrt{100}$) y llega a la longitud exacta de $10$. \\
\textbf{Veredicto:} \textbf{Correcto} (2/2 pts).

\subsection*{3. Vectores (Fuerza resultante)}
\textbf{Procedimiento y resultado:} Agrupa las componentes en los ejes correspondientes $\hat{i}$ y $\hat{j}$. Sin embargo, comete un \textbf{error de signo} en la componente $\hat{j}$: en lugar de sumar $(25 + 35)\hat{j}$, realiza una resta $(25 - 35)\hat{j}$. Esto la lleva a un vector incorrecto de $30\hat{i} - 10\hat{j}$. \\
\textbf{Veredicto:} \textbf{Parcialmente correcto} (Planteamiento válido, error de signo al operar) (1/2 pts).

\subsection*{4. Potencias de diez (Notación científica y prefijo)}
\textbf{Procedimiento y resultado:} Transforma la cifra correctamente a notación científica ($1,25 \times 10^{-5}$) y la ajusta de forma excelente a $12,5 \times 10^{-6}$ para asociarla a un prefijo. Sin embargo, omite la respuesta final, que consistía en identificar y escribir el nombre o símbolo de dicho prefijo (micrómetros, $\mu$m). \\
\textbf{Veredicto:} \textbf{Incompleto / Parcialmente correcto} (1/2 pts).

\subsection*{5. Sistemas de unidades (Conversión de densidad)}
\textbf{Procedimiento y resultado:} El primer planteamiento formal de factores de conversión en la parte superior es incorrecto. Afortunadamente, en la parte inferior desglosa las operaciones aritméticas paso a paso de manera perfecta: divide entre 1000 y luego multiplica por 1.000.000, llegando al resultado exacto de $1030$. \\
\textbf{Veredicto:} \textbf{Correcto} (2/2 pts).

\vspace{0.5cm}
\section*{Comentarios Generales y Sugerencias}
La estudiante demuestra un buen dominio en el cálculo aritmético básico y la notación de potencias. Se recomiendan los siguientes puntos de mejora:
\begin{itemize}
    \item \textbf{Cuidado con los signos:} Es fundamental repasar la suma algebraica de vectores para evitar restar los componentes cuando el enunciado solicita la fuerza resultante (Pregunta 3).
    \item \textbf{Revisión de respuestas:} Prestar mayor atención para no dejar el ejercicio incompleto (como el nombre del prefijo en la Pregunta 4).
\end{itemize}

\end{document}